\section*{Chapter \ref{ch:b2:wilson-cycle} --- Ancient Massifs and the Wilson Cycle}

\begin{solution}{exo:b2:wilson-cycle:1}
(a) $600 - 400 = \qty{200}{Myr}$, from the rifting of the margins to the last collision.
(b) $480 - 400 = \qty{80}{Myr}$. (c) The Caledonian orogeny is a set of phases spread over
some 200 million years, and each phase began at different times in different places: the word
locates an ocean and its belt, not a moment. An age is given by the stratigraphic scale or in
millions of years.
\end{solution}

\begin{solution}{exo:b2:wilson-cycle:2}
(c) the East African rift, rifting; (a) the Red Sea, a young ocean; (d) the Atlantic, a wide
ocean with passive margins; (e) the Pacific margin of South America, subduction and the
closing of an ocean; (b) the Himalaya, collision; (f) Brittany, a belt worn flat, after which
rifting may start again.
\end{solution}

\begin{solution}{exo:b2:wilson-cycle:3}
(a) After the deposition of the early Silurian beds, so after \qty{443.1}{Ma}, and before the
deposition of the Late Devonian beds, so before \qty{358.9}{Ma}, the end of the Devonian.
(b) Less than $443.1 - 358.9 = \qty{84.2}{Myr}$; in fact less, since the lower beds are younger
than the base of the Silurian and the upper beds older than its top. (c) Deposition of the
grey turbidites in the ocean; folding and cleavage in the Caledonian collisions; uplift and
erosion to an uneven land surface; deposition of the red breccias and conglomerates; gentle
tilting in the Variscan orogeny; today's erosion, which cuts the cliff.
\end{solution}

\begin{solution}{exo:b2:wilson-cycle:4}
(a) Both: a block of continental crust separated from the continents by oceans while it
crossed the Iapetus, and, once docked, a fault-bounded block whose history differs from its
neighbours'. (b) Neither: a continent, not a block added to one. (c) A \omterm{def:b2:wilson-cycle:terrane}{terrane}, not a
\omterm{def:b2:wilson-cycle:terrane}{microcontinent}: its crust is that of an arc, not a fragment of a continent. (d) Neither: it is
the boundary between \omterm{def:b2:wilson-cycle:terrane}{terranes}, the trace of the closed ocean.
\end{solution}

\begin{solution}{exo:b2:wilson-cycle:5}
(a) $3000 - 2600 = 400$, $2600 - 1500 = 1100$, $1500 - 875 = 625$, $875 - 600 = 275$, $600 -
262.5 = 337.5$; mean $2737.5/5 = \qty{547.5}{Myr}$. (b) $1100/275 = 4.0$. (c) $(1500 -
262.5)/3 = \qty{412.5}{Myr}$. The mean moves by a quarter when the doubtful points are dropped
and the spacings vary fourfold: a rough recurrence, not a period.
\end{solution}

\begin{solution}{exo:b2:wilson-cycle:6}
(a) $3000/(70 - 25) = \qty{66.7}{Myr}$. (b) $3000/70 = \qty{42.9}{Myr}$. (c) Once the ridge has
been subducted, the ocean floor left is part of the far plate, which carries the far
continent: the closing rate is the velocity of that plate relative to the trench.
\end{solution}

\begin{solution}{exo:b2:wilson-cycle:7}
(a) $700/8 = \qty{87.5}{km}$; $700/30 = \qty{23.3}{km}$. (b) The Variscan belt: at
\qty{30}{\celsius} per kilometre the middle crust melts; at \qty{8}{\celsius} per kilometre
the melting temperature lies below the base of even a thickened crust. (c) The low gradient
describes rocks carried down fast in a subduction zone, before they could heat up. Rocks that
stay deep in a thickened crust warm towards a normal geotherm over millions of years, and can
then reach melting at smaller depth.
\end{solution}

\begin{solution}{exo:b2:wilson-cycle:8}
(a) $\tau\ln(3/0.3) = 26.3\times2.303 = \qty{60.6}{Myr}$. (b) $290 - 60.6 = \qty{229}{Ma}$, in
the Late Triassic. (c) $\tau = (290 - 251.9)/\ln10 = \qty{16.5}{Myr}$. (d) The extensional
collapse of the thickened crust, which thins it and lowers its surface directly; also a
higher erosion rate, a wetter climate or softer rocks, or a lower starting elevation.
\end{solution}

\begin{solution}{exo:b2:wilson-cycle:9}
(a) The suture, the trace of the Iapetus, lies on land in Britain and Ireland, east of the
Atlantic; crust of Laurentia, from the far side of the Iapetus, lies near Rockall on the
European side. The new rift cut west of the old suture: it opened near it, not along it.
(b) North of the suture, basement and cover of the Laurentian margin; south of it, those of
Avalonia, with different histories and, in the early Palaeozoic, different faunas. (c) Rifting
came after 80 million years, more than the 30 to 50 after which a cooled zone may have grown
strong; yet the Central Atlantic did follow the Appalachians. The heat of a belt fades; its
faults, fabrics and compositions remain and can still guide a rift.
\end{solution}

\begin{solution}{exo:b2:wilson-cycle:10}
(a) $70 - 40/2 = \qty{50}{mm}$ a year. (b) $\qty{1500}{km}/\qty{50}{km.Myr^{-1}} =
\qty{30}{Myr}$. (c) The ocean narrows at $70 - 40 = \qty{30}{mm}$ a year: $3000 - 30\times30 =
\qty{2100}{km}$. (d) Young, hot lithosphere then arrives at the trench, and no old, cold slab
dehydrates beneath the arc any more: the arc magmatism changes or stops.
\end{solution}

\begin{solution}{exo:b2:wilson-cycle:11}
(a) $h = \Delta T(\rho_m - \rho_c)/\rho_m = 25\times600/3300 = \qty{4.5}{km}$. (b) All the extra
\qty{25}{km}, 5.5 times the height of the range. (c) $2700\times9.8\times25\,000 =
\qty{6.6e8}{Pa}$, about \qty{0.66}{GPa}. (d) At about \qty{0.7}{GPa} and the temperatures of
the root of a thickened crust, several hundred degrees, the amphibolite or the granulite
facies: the gneisses and \omterm{def:b2:crustal-melting:migmatite}{migmatites} of the old massifs.
\end{solution}

\begin{solution}{exo:b2:wilson-cycle:12}
(a) $S = 30\times10 = \qty{300}{km}$; $W_0 = 300 + 300 = \qty{600}{km}$; $T = 35\times600/300 =
\qty{70}{km}$. (b) $h = 35\times600/3300 = \qty{6.36}{km}$; $r = 4.5h = \qty{28.6}{km}$.
(c) $\Delta F = 2700\times9.8\times6364\times(35\,000 + 3182 + 14\,318) =
\qty{8.8e12}{N.m^{-1}}$, 1.4 times Tibet's. (d) Such a belt is heavier than Tibet, which already
spreads: \omterm{def:b2:himalaya-andes:collapse}{gravitational collapse}, by extension of its high regions and thrusting at its edges.
\end{solution}

\begin{solution}{pb:b2:wilson-cycle:1}
\textbf{1.} $720/(540 - 500) = \qty{18}{km}$ per million years, \qty{18}{mm} a year.
\textbf{2.} \qty{36}{mm} a year: a slow-spreading ridge (20 to \qty{55}{mm} a year).
\textbf{3.} $36\times(580 - 540) = \qty{1440}{km}$.
\textbf{4.} $36\times(580 - 500) = \qty{2880}{km}$.
\textbf{5.} Symmetric spreading at a constant rate from zero width at \qty{580}{Ma}, and no
subduction before \qty{500}{Ma}.
\textbf{6.} \qty{80}{Myr}, at both edges of the ocean, against the continental margins. It is the
coldest, thickest and densest oceanic lithosphere of the ocean, the most ready to sink.
\textbf{7.} $64 - 36 = \qty{28}{mm}$ a year.
\textbf{8.} $2880/28 = \qty{102.9}{Myr}$ after \qty{500}{Ma}: at \qty{397.1}{Ma}.
\textbf{9.} The ridge, \qty{1440}{km} from the trench, approaches it at $64 - 18 =
\qty{46}{mm}$ a year: after $1440/46 = \qty{31.3}{Myr}$, at \qty{468.7}{Ma}.
\textbf{10.} $2880 - 28\times31.3 = \qty{2004}{km}$.
\textbf{11.} $2004/40 = \qty{50.1}{Myr}$ later: at $468.7 - 50.1 = \qty{418.6}{Ma}$.
\textbf{12.} Inside the 425 to \qty{400}{Ma} of the Scandian collision.
\textbf{13.} $40\times8 = \qty{320}{km}$.
\textbf{14.} $W_0 = 400 + 320 = \qty{720}{km}$; $T = 35\times720/400 = \qty{63.0}{km}$.
\textbf{15.} $h = 28\times600/3300 = \qty{5.09}{km}$; $r = 4.5\times5.09 = \qty{22.9}{km}$.
\textbf{16.} $\Delta F = 2700\times9.8\times5091\times(35\,000 + 2545 + 11\,455) =
\qty{6.60e12}{N.m^{-1}}$, about \qty{8}{\percent} more than Tibet's.
\textbf{17.} To stop rising and spread: extension and normal faults in its high interior,
thrusts at its edges, as Tibet does today.
\textbf{18.} $f = 600/3300 = 0.182$; $k = 2.303/11 = 0.209$ per million years; $\tau = 1/(kf)
= \qty{26.3}{Myr}$.
\textbf{19.} $3.0\,e^{-30/26.3} = \qty{0.96}{km}$.
\textbf{20.} $26.3\ln(3.0/0.5) = \qty{47.1}{Myr}$ after \qty{410}{Ma}: at \qty{362.9}{Ma}.
\textbf{21.} The extra crust is $h/f$: $3.0/0.182 = \qty{16.5}{km}$ falls to $0.5/0.182 =
\qty{2.75}{km}$, so \qty{13.75}{km} of rock is removed.
\textbf{22.} $2700\times9.8\times13\,750 = \qty{3.6e8}{Pa}$, about \qty{0.36}{GPa}.
\textbf{23.} Yes: by the Late Devonian the model belt is down to about \qty{1}{km} of mean
elevation and still falling, with relief left between its valleys and ridges; red gravels can
gather in its lows well before \qty{358.9}{Ma}, as at Siccar Point, where they rest on a surface
with relief.
\textbf{24.} \textbf{\qty[text-series-to-math=true]{2880}{km}}, at \qty{500}{Ma}, when
subduction began.
\end{solution}
